\documentclass[12pt]{article}
\usepackage[utf8]{inputenc}
\usepackage[T1]{fontenc}
\usepackage{hyperref}
\title{Poumai Naga Language Resources for K-12 Curriculum Development}
\author{Compiled by Claude Sonnet 5.5 (claude-sonnet-5-5) for Andr\'as Kornai}
\date{}
\begin{document}
\maketitle
\section{Introduction}
Poumai Naga (Poula; ISO 639-3 pmx) is a Sino-Tibetan (Trans-Himalayan) language spoken mainly in Manipur, north-east India. Little published material exists online; this report lists only what turned up in web searches. No school textbooks or curricula in the language were found.
\section{Primary Educational Resources}
No primers, textbooks or school programmes for Poumai Naga were found in searches.
\section{Linguistic and Reference Materials}
\begin{itemize}
\item Grammar of Poumai Naga (Poula): a Trans-Himalayan language of north-east India (S. L. Veikho, Brill, 2021; Bern thesis 2019), with a bilingual lexicon and analysed texts in the appendices. Catalogue record: \url{https://ci.nii.ac.jp/ncid/BC11892297}
\item Poula language, encyclopedia overview (tonal; /\textschwa/ written with \"u or ii in Bible and hymn orthography): \url{https://en.wikipedia.org/wiki/Poula_language}
\end{itemize}
A Poula Dialects Dictionary by the Poumai Literature Society is reported as in progress; no published edition was confirmed.
\section{Religious and Community Resources}
The orthography used in Bible and hymn texts was devised by the Bible Society of India (per the overview above). No Bible edition URL was found in searches.
\section{Digital and Audio Resources}
\begin{itemize}
\item Poumai Naga Language Resource collection, University of North Texas Digital Library: \url{https://digital.library.unt.edu/explore/collections/PNLR/}
\end{itemize}
\section{Curriculum Development Context}
Documentation is mostly descriptive-linguistic. A curriculum would need to build on the Veikho grammar and lexicon and on community bodies such as the Poumai Literature Society; tone is not marked in the current orthography.
\end{document}